\begin{table*}[!t]
\caption{Retrieval recovery (Panel A): 100 CategoryBqueries,342 gold authorities.}
\label{tab:retrieval_eval}
\centering\begingroup\fontsize{8}{9.6}\selectfont\setlength{\tabcolsep}{3pt}\renewcommand{\arraystretch}{1.15}
\begin{tabular}{lcccc}
\toprule
Pipeline & Macro recall & Micro recall & Complete targets & Count \\
\midrule
Flat RAG & \CatBFlatMacroNodeRecall & \CatBFlatMicroNodeRecall & \CatBFlatFullPathRecall & \CatBFlatFullPathCount \\
GraphRAG & \CatBGraphRAGMacroNodeRecall & \CatBGraphRAGMicroNodeRecall & \CatBGraphRAGFullPathRecall & \CatBGraphRAGFullPathCount \\
HippoRAG & \CatBHippoRAGMacroNodeRecall & \CatBHippoRAGMicroNodeRecall & \CatBHippoRAGFullPathRecall & \CatBHippoRAGFullPathCount \\
TERRA & \CatBTerraMacroNodeRecall & \CatBTerraMicroNodeRecall & \CatBTerraFullPathRecall & \CatBTerraFullPathCount \\
\bottomrule
\end{tabular}
\endgroup\par\vspace{2pt}
\raggedright\footnotesize Panel A above uses canonical stored contexts; target-set matching does not certify connected paths. Panel B chain-length cells and Panel C all 22 sensitivity configurations are retained in the corresponding full-detail supplement. Canonical 66/100, legacy 62/100 and bounded 61/100 are distinct. Five-node stratum n 5 is descriptive. Full cells, methods and notes: Supplement Section~\ref{supp-detail-tab:retrieval_eval}.
\end{table*}
